\documentclass[10pt,landscape]{article}

\usepackage[margin=0.32in]{geometry}
\usepackage[T1]{fontenc}
\usepackage[utf8]{inputenc}
\usepackage{lmodern}
\usepackage{textcomp}
\usepackage{microtype}
\usepackage[table,dvipsnames]{xcolor}
\usepackage{array}
\usepackage{booktabs}
\usepackage{enumitem}
\usepackage{hyperref}
\usepackage{listings}
\usepackage{multicol}
\usepackage{tabularx}
\usepackage{titlesec}

\pagestyle{empty}
\setlength{\parindent}{0pt}
\setlength{\parskip}{1.5pt}
\setlength{\columnsep}{13pt}
\setlength{\columnseprule}{0.2pt}
\setcounter{secnumdepth}{0}

\definecolor{ToshBlue}{HTML}{0F4C81}
\definecolor{ToshInk}{HTML}{1F2937}
\definecolor{ToshMuted}{HTML}{5B6470}
\definecolor{ToshPanel}{HTML}{F4F7FB}
\definecolor{ToshRule}{HTML}{CBD5E1}
\definecolor{ToshCodeBg}{HTML}{F8FAFC}

\color{ToshInk}

\titleformat{\section}
  {\large\bfseries\color{ToshBlue}}
  {}
  {0pt}
  {}
\titlespacing*{\section}{0pt}{0.4em}{0.12em}

\setlist[itemize]{leftmargin=1.05em,itemsep=0.5pt,topsep=1.5pt,parsep=0pt,partopsep=0pt}

\lstdefinestyle{tosh}{
  basicstyle=\ttfamily\scriptsize,
  backgroundcolor=\color{ToshCodeBg},
  frame=single,
  rulecolor=\color{ToshRule},
  framerule=0.4pt,
  xleftmargin=0pt,
  xrightmargin=0pt,
  aboveskip=0.18em,
  belowskip=0.24em,
  columns=fullflexible,
  keepspaces=true,
  showstringspaces=false,
  breaklines=true,
  literate={°}{{\textdegree}}1
           {μ}{{\ensuremath{\mu}}}1
           {Ω}{{\ensuremath{\Omega}}}1
}
\lstset{style=tosh}

\newcommand{\sheettitle}[3]{
  {\LARGE\bfseries\color{ToshBlue} ToSh Cheat Sheet: #1\par}
  \vspace{0.08em}
  {\small\color{ToshMuted} #2 \hfill #3\par}
  \vspace{0.18em}
  {\color{ToshRule}\rule{\linewidth}{0.8pt}}
  \vspace{0.12em}
}

\newcommand{\note}[1]{
  \vspace{0.12em}
  \fcolorbox{ToshRule}{ToshPanel}{%
    \parbox{\dimexpr\linewidth-2\fboxsep-2\fboxrule\relax}{\footnotesize #1}%
  }
  \vspace{0.12em}
}

\newcommand{\tighttable}[2]{
  {\footnotesize
  \renewcommand{\arraystretch}{1.08}
  \begin{tabularx}{\linewidth}{@{}#1@{}}
  #2
  \end{tabularx}}
}


\begin{document}
\footnotesize
\sheettitle{Shell Basics}{REPL workflow, help, capture forms, redirection, and jobs}{spec + parser/tests}

\begin{multicols*}{3}

\section{Start Here}
\tighttable{>{\ttfamily\raggedright\arraybackslash}p{0.54\linewidth}X}{
\toprule
Command & What it does \\
\midrule
dotnet run --project src/Tosh.Cli & Start the ToSh REPL. \\
help ls & Explain one command or topic. \\
help search regex & Search the help catalog. \\
help browse grep & Open the full-screen help browser. \\
time \{ ls -la \} & Measure wall time, CPU, and memory. \\
history search ls & Search command history by text. \\
\bottomrule
}

\begin{lstlisting}
# typed object pipeline
ls -la | where _.Type == file | sort Size | reverse | first 5

# shell-friendly discovery
help which
apropos json
\end{lstlisting}

\section{Mental Model}
\begin{itemize}
\item Pipelines carry real CLR objects, not plain text by default.
\item The current pipeline item is \texttt{\_} inside \texttt{where}, \texttt{each}, \texttt{map}, and similar blocks.
\item \texttt{\$tosh.Last.Result} stores the last successful statement result.
\end{itemize}

\note{\textbf{Rule of thumb:} text is something you opt into with commands like \texttt{raw}, \texttt{to}, \texttt{join-lines}, \texttt{\$(...)} or external programs.}

\section{Capture Forms}
\tighttable{>{\ttfamily\raggedright\arraybackslash}p{0.28\linewidth}X}{
\toprule
Form & Meaning \\
\midrule
(...) & Evaluate a pipeline/expression and require exactly one object value. \\
\$(...) & Capture text output; multiple values join with newlines. \\
<(...) & Materialize a pipeline to a temp file and pass its path-like handle. \\
>(...) & Send output to a process-substitution sink. \\
\bottomrule
}

\begin{lstlisting}
var user = (whoami)
var text = $(echo one two three)
var names = <(ls | get Name | sort)
/bin/cat $names
\end{lstlisting}

\note{\texttt{(...)} is strict: if the inner pipeline yields 0 or more than 1 value, ToSh raises a diagnostic.}

\section{Session Commands}
\tighttable{>{\ttfamily\raggedright\arraybackslash}p{0.34\linewidth}X}{
\toprule
Command & Notes \\
\midrule
pwd, cd, dirs & Current directory and stack navigation. \\
back, forward & Move through the directory stack. \\
env, export, unset & Read or change environment/session names. \\
vars & Show visible variables in scope. \\
which ls & Resolve a command name. \\
clear & Clear the terminal display. \\
\bottomrule
}

\begin{lstlisting}
pwd
cd src
dirs
back
env USER
\end{lstlisting}

\columnbreak

\section{Redirection}
\tighttable{>{\ttfamily\raggedright\arraybackslash}p{0.27\linewidth}X}{
\toprule
Syntax & Meaning \\
\midrule
out> file & stdout to file, truncate \\
out>> file & stdout to file, append \\
err> file & stderr to file, truncate \\
o+e> file & stdout and stderr together \\
<<< text & here-string input \\
\bottomrule
}

\begin{lstlisting}
echo "hello" out> hello.txt
echo "world" out>> hello.txt
/bin/sh -c "echo boom >&2" err> errors.txt
grep toast <<< "toast shell"
\end{lstlisting}

\note{ToSh uses named redirection because \texttt{<} and \texttt{>} are also comparison operators in expressions.}

\section{Jobs And Async Work}
\begin{lstlisting}
/bin/sleep 5 &
jobs
wait-for 1
kill 1
\end{lstlisting}

\begin{itemize}
\item A trailing \texttt{\&} backgrounds a pipeline or external command.
\item \texttt{jobs} returns typed job records you can filter or inspect.
\item \texttt{wait-for} waits on job ids or piped jobs.
\end{itemize}

\note{At the \emph{start} of an expression, \texttt{\&name} means ``function reference'', not ``background job''.}

\section{Useful One-Liners}
\begin{lstlisting}
ls -la | where _.Type == file | get { Name, Size }
ip addr | where { _.State == up }
df | summarize _.Used
\end{lstlisting}

\begin{lstlisting}
guid new v7
time { ls -la }
\end{lstlisting}

\end{multicols*}
\newpage

\sheettitle{Shell Basics}{Interactive browsers, history, prompt modules, and live config}{page 2 / front-back reference}

\begin{multicols*}{3}

\section{Help And Discovery}
\begin{lstlisting}
help --cli regex
help related ls
help categories
apropos pipeline
which ls grep tosh
types list
\end{lstlisting}

\begin{itemize}
\item \texttt{help --cli} opens the inline fuzzy browser.
\item \texttt{help browse} opens the full-screen split-pane browser.
\item \texttt{types} searches both CLR types and ToSh shell types.
\end{itemize}

\section{Config And Display}
\begin{lstlisting}
config browse prompt
config get prompt.name-text
config set prompt.name-text toast
config reset display
config reload
\end{lstlisting}

\begin{lstlisting}
view detail
view compact
view columns
view size bytes
view datetime table unix
\end{lstlisting}

\note{\texttt{config} manages live shell settings and startup config; \texttt{view} is the quick way to tune rendering for the current session.}

\section{History Workflow}
\begin{lstlisting}
history
history path
history search grep
history expand 237
history run 237
history save
\end{lstlisting}

\begin{itemize}
\item REPL sugar: \texttt{!!}, \texttt{!237}, \texttt{!-2}, \texttt{!?text?}, \texttt{\^{}old\^{}new\^{}}
\item \texttt{history delete <spec>} removes entries; \texttt{reload} re-reads the file-backed store.
\end{itemize}

\section{Directory Stack}
\begin{lstlisting}
dirs
dirs goto 2
dirs remove 1
dirs clear
\end{lstlisting}

\begin{itemize}
\item \texttt{dirs} is handy when you bounce among a few working roots.
\item Use \texttt{goto} for fast return without retyping long paths.
\end{itemize}

\columnbreak

\section{Inspecting Objects}
\begin{lstlisting}
ls -la | first | inspect
new System.Random() | inspect -a
new System.Random() | inspect --flat
name-of (pwd)
type-of (date now)
\end{lstlisting}

\begin{lstlisting}
ls -la | first | get-props
ls -la | first | get-prop Name
has-prop (pwd) FullName
has-method (pwd) ToString
\end{lstlisting}

\note{In the REPL, \texttt{F1} opens help for the token under the cursor, \texttt{F2} tries inline inspect, and \texttt{Alt+H} / \texttt{Alt+I} are fallback bindings.}

\section{Prompt Modules}
\tighttable{>{\ttfamily\raggedright\arraybackslash}p{0.34\linewidth}X}{
\toprule
Module & What it adds \\
\midrule
\texttt{prompt-dir} & current directory \\
\texttt{prompt-git} & branch and repo status \\
\texttt{prompt-time} & current wall clock time \\
\texttt{prompt-userhost} & user and host name \\
\texttt{prompt-history} & next history id \\
\texttt{prompt-jobs} & current background job count \\
\texttt{prompt-text} & literal styled text \\
\texttt{prompt-newline} & multi-line prompt split \\
\bottomrule
}

\begin{lstlisting}
prompt-userhost --dim
prompt-dir --depth 2 --fg blue
prompt-git --fg bright-green
prompt-time --format "HH:mm"
\end{lstlisting}

\section{Run Files And Processes}
\begin{lstlisting}
source ./script.tosh
source ./script.tosh ...$tosh.Script.Args
env DEBUG=1 -- dotnet test
exec zsh
exit
\end{lstlisting}

\section{Jobs And Waiting}
\begin{lstlisting}
sleep 5s &
jobs
wait-for $tosh.Last.Result.Id
kill 3
\end{lstlisting}

\begin{itemize}
\item \texttt{jobs} and \texttt{wait-for} are the main background-job tools.
\item \texttt{bg} and \texttt{fg} resume suspended jobs stopped with \texttt{Ctrl+Z}.
\end{itemize}

\columnbreak

\section{Runtime State}
\begin{lstlisting}
vars
env USER HOME
$tosh.Last.Result
$tosh.Script.Path
$tosh.Session.OpenHandleCount
\end{lstlisting}

\begin{itemize}
\item \texttt{\$tosh.Last.Result} tracks the last successful statement result.
\item \texttt{\$tosh.Script.Path} is useful inside sourced or executed scripts.
\item Open managed handles are visible through the \texttt{\$tosh.Session} namespace.
\end{itemize}

\section{Small Conveniences}
\begin{lstlisting}
assert ($env.HOME != null)
styled "warning" --fg yellow --bold
read-line "name> "
ignore
\end{lstlisting}

\tighttable{>{\ttfamily\raggedright\arraybackslash}p{0.31\linewidth}X}{
\toprule
Tool & Why it matters \\
\midrule
\texttt{assert} & fail early in scripts with a clear message \\
\texttt{styled} & emit colored/styled terminal text segments \\
\texttt{read-line} & prompt for one line of interactive text \\
\texttt{ignore} & explicitly discard pipeline output \\
\bottomrule
}

\section{Config Object}
\begin{lstlisting}
$tosh.Config.Prompt.NameText = "toast"
$tosh.Config.Repl.ContinuationPrompt = "..> "
$tosh.Config.Display.StorageSize.Mode = "Bytes"
config get prompt
\end{lstlisting}

\begin{itemize}
\item Direct \texttt{\$tosh.Config} edits affect the live session immediately.
\item Use \texttt{config reload} to replay startup config files into the session.
\end{itemize}

\note{\texttt{source} reruns a file every time in the current scope; \texttt{require} caches modules once per session; \texttt{exec} replaces the current ToSh process with another program.}

\end{multicols*}
\end{document}
